\appendix

\chapter{Tables}
\label{app:tables}

All formulas use the convention
$\F[f](\omega)=\int f(t)e^{-i\omega t}\dd t$ with $2\pi$ on the inverse transform, and pairings
without complex conjugation.

\section{Transform pairs}

\begin{center}
\small
\begin{tabular}{@{}lll@{}}
\toprule
Time domain & Frequency domain & Where proved \\
\midrule
$\delta(t)$ & $1$ & \eqref{eq:Fdelta} \\
$\delta^{(n)}(t)$ & $(i\omega)^{n}$ & \eqref{eq:Fdelta} \\
$1$ & $2\pi\delta(\omega)$ & \eqref{eq:Fdelta} \\
$t^{\,n}$ & $2\pi i^{\,n}\delta^{(n)}(\omega)$ & \eqref{eq:Ftn} \\
$\delta(t-a)$ & $e^{-ia\omega}$ & \eqref{eq:timeshift} \\
$e^{i\omega_{0}t}$ & $2\pi\delta(\omega-\omega_{0})$ & \eqref{eq:modulation} \\
$\cos(\omega_{0}t)$ & $\pi\big[\delta(\omega-\omega_{0})+\delta(\omega+\omega_{0})\big]$ & linearity \\
$\sin(\omega_{0}t)$ & $(\pi/i)\big[\delta(\omega-\omega_{0})-\delta(\omega+\omega_{0})\big]$ & linearity \\
$H(t)$ & $\pi\delta(\omega)-i\,\PV(1/\omega)$ & \eqref{eq:FHgen} \\
$\sgn t$ & $-2i\,\PV(1/\omega)$ & \eqref{eq:sgn} \\
$\PV(1/t)$ & $-i\pi\,\sgn\omega$ & duality, \eqref{eq:F2} \\
$e^{-\ep|t|}$ & $2\ep/(\ep^{2}+\omega^{2})$ & Theorem~\ref{thm:exptransform} \\
$H(t)e^{-\ep t}$ & $1/(\ep+i\omega)$ & Theorem~\ref{thm:Heps} \\
$e^{-at^{2}/2}$, $a>0$ & $\sqrt{2\pi/a}\;e^{-\omega^{2}/(2a)}$ & Lemma~\ref{lem:gauss} \\
\bottomrule
\end{tabular}
\end{center}

\begin{remark}[the sine entry]
Write $\sin(\omega_{0}t)=\frac{1}{2i}(e^{i\omega_{0}t}-e^{-i\omega_{0}t})$ and transform term by
term. The coefficient is $\pi$, not $2\pi$, because sine is half the difference of two exponentials.
Some tables list $\pi i[\delta(\omega-\omega_{0})-\delta(\omega+\omega_{0})]$, which corresponds to
the opposite sign convention in the exponent; the two differ by a sign, and neither oddness nor
vanishing at the origin distinguishes them, so the derivation must be redone rather than recalled.
\end{remark}

\section{Operational rules}

\begin{center}
\small
\begin{tabular}{@{}lll@{}}
\toprule
Rule & Formula & Hypothesis \\
\midrule
time shift & $\F[f(\cdot-a)]=e^{-ia\omega}\fh$ & \eqref{eq:timeshift} \\
modulation & $\F[e^{i\omega_{0}t}f]=\fh(\omega-\omega_{0})$ & \eqref{eq:modulation} \\
dilation & $\F[f(a\,\cdot)]=|a|^{-1}\fh(\omega/a)$ & $a\neq0$ \\
multiply by $t$ & $\F[tf]=i\,\partial_{\omega}\fh$ & $f\in\Sw$ \\
differentiate & $\F[f']=i\omega\fh$ & $f\in\Sw$ \\
distributional & $\F(tT)=i\,\partial_{\omega}\widehat T$ & \eqref{eq:multt-dist} \\
distributional & $\F(T')=i\omega\widehat T$ & \eqref{eq:deriv-dist} \\
twice & $\F^{2}T=2\pi\,T(-\cdot)$ & \eqref{eq:F2} \\
inverse & $\F^{-1}T=(2\pi)^{-1}\F(T(-\cdot))$ & Theorem~\ref{thm:inverse} \\
\midrule
scale $\delta$ & $\delta(ax)=|a|^{-1}\delta(x)$ & $a\neq0$ \\
compose $\delta$ & $\delta(g(x))=\sum_{j}\delta(x-x_{j})/|g'(x_{j})|$ & $g$ has simple zeros \\
multiply $\delta$ & $a\,\delta=a(0)\,\delta$ & $a\in C^{\infty}$ \\
multiply by $x$ & $x\,\delta^{(n)}=-n\,\delta^{(n-1)}$ & \eqref{eq:xdeltaN} \\
Leibniz & $(aT)'=a'T+aT'$ & $a\in C^{\infty}$ \\
jumps & $Df=f'_{\mathrm{cl}}+\sum_{j}J_{j}\delta_{a_{j}}$ & $f$ piecewise $C^{1}$ \\
\bottomrule
\end{tabular}
\end{center}

\section{Integrals used}

\begin{center}
\small
\begin{tabular}{@{}lll@{}}
\toprule
Integral & Value & Used in \\
\midrule
$\displaystyle\int_{0}^{\infty}e^{-at}\dd t$ & $1/a$ \ ($\Re a>0$) & Theorem~\ref{thm:exptransform} \\
$\displaystyle\int_{\R}\frac{\dd\omega}{\ep^{2}+\omega^{2}}$ & $\pi/\ep$ & Proposition~\ref{prop:three} \\
$\displaystyle\int_{\R}\frac{\ep\dd\omega}{\ep^{2}+\omega^{2}}$ & $\pi$ & \eqref{eq:masspi} \\
$\displaystyle\int_{\R}\frac{2\ep\dd\omega}{\ep^{2}+\omega^{2}}$ & $2\pi$ & \eqref{eq:three} \\
$\displaystyle\int_{-a}^{a}\frac{\ep\dd\omega}{\pi(\ep^{2}+\omega^{2})}$ &
$\frac{2}{\pi}\arctan\frac{a}{\ep}$ & \eqref{eq:concentration} \\
$\displaystyle\int_{\R}e^{-u^{2}/2}\dd u$ & $\sqrt{2\pi}$ & Lemma~\ref{lem:gauss} \\
$\displaystyle\int_{\R}e^{-\eta\omega^{2}/2}e^{i\omega x}\dd\omega$ &
$\sqrt{2\pi/\eta}\,e^{-x^{2}/(2\eta)}$ & \eqref{eq:gaussint} \\
$\displaystyle\int_{-T}^{T}e^{-i\omega t}\dd t$ & $2\sin(\omega T)/\omega$ & \eqref{eq:IT} \\
$\displaystyle\int_{\R}\frac{2\ep x}{\pi(x^{2}+\ep^{2})^{2}}\dd x$ & $1/(\pi\ep)$ & Proposition~\ref{prop:Kprime} \\
\bottomrule
\end{tabular}
\end{center}

\section{Which space does an object belong to}

\begin{center}
\small
\begin{tabular}{@{}lllll@{}}
\toprule
Object & $\Lone$ & $\Ltwo$ & $\D'$ & $\Sw'$ \\
\midrule
$e^{-\ep|t|}$ & yes & yes & yes & yes \\
$1$ & no & no & yes & yes \\
$t^{\,n}$ & no & no & yes & yes \\
$H$ & no & no & yes & yes \\
$\delta$, $\delta^{(n)}$ & no & no & yes & yes \\
$\PV(1/\omega)$ & no & no & yes & yes \\
$1/(1+t^{2})$ & yes & yes & yes & no \\
$e^{t^{2}}$ & no & no & yes & no \\
\bottomrule
\end{tabular}
\end{center}

\begin{remark}
Two rows carry most of the information. The constant $1$ is in neither $\Lone$ nor $\Ltwo$, which is
why neither the elementary theory nor Plancherel can be applied to it; it is in $\Sw'$, which is why
the generalized transform can. The function $e^{t^{2}}$ is locally integrable and therefore defines a
distribution, but it is not tempered and consequently has no Fourier transform in this framework.
Membership is a property of the space, not of the formula.
\end{remark}

\section{Conventions in use}

\begin{center}
\small
\begin{tabular}{@{}llll@{}}
\toprule
Convention & Forward transform & $\F[1]$ & Relation to this text \\
\midrule
angular (here) & $\displaystyle\int f(t)e^{-i\omega t}\dd t$ & $2\pi\delta(\omega)$ & --- \\
symmetric & $\displaystyle\frac{1}{\sqrt{2\pi}}\int f(t)e^{-i\omega t}\dd t$ &
$\sqrt{2\pi}\,\delta(\omega)$ & multiply by $\frac{1}{\sqrt{2\pi}}$ \\
frequency & $\displaystyle\int f(t)e^{-2\pi i\nu t}\dd t$ & $\delta(\nu)$ & set $\omega=2\pi\nu$ \\
\bottomrule
\end{tabular}
\end{center}

The conversion rule is
\begin{equation}
\F_{\nu}[f](\nu)=\F_{\omega}[f](2\pi\nu),
\qquad
\delta(2\pi\nu)=\frac{1}{2\pi}\delta(\nu).
\end{equation}

\clearpage
\section{Where each statement is proved}

\begin{center}
\small
\begin{tabular}{@{}lll@{}}
\toprule
Statement & First obtained & Fully justified \\
\midrule
$\F[e^{-\ep|\cdot|}]=2\ep/(\ep^{2}+\omega^{2})$ & Theorem~\ref{thm:exptransform} & same (classical) \\
$\int K_{\ep}\varphi\to\varphi(0)$ & Theorem~\ref{thm:sampling} & Theorem~\ref{thm:sampling-general} \\
$\F[1]=2\pi\delta$ & \eqref{eq:F1} & Theorem~\ref{thm:main} \\
$\ip{\delta^{(n)}}{\varphi}=(-1)^{n}\varphi^{(n)}(0)$ & \eqref{eq:deltan} & \eqref{eq:xdeltaN}--\eqref{eq:leibniz} \\
$\F[t^{\,n}]=2\pi i^{\,n}\delta^{(n)}$ & Theorem~\ref{thm:tn} & Theorem~\ref{thm:main} \\
$\F[H]=\pi\delta+\PV(1/i\omega)$ & \eqref{eq:FH} & Theorem~\ref{thm:main} \\
$\F[\sgn t]=-2i\PV(1/\omega)$ & \eqref{eq:sgn} & Theorem~\ref{thm:main} \\
$\F[\PV(1/t)]=-i\pi\sgn\omega$ & Proposition~\ref{prop:PVtransform} & Theorem~\ref{thm:F2} \\
inversion $\varphi=\frac{1}{2\pi}\int\ph e^{i\omega\cdot}$ & Theorem~\ref{thm:inversion} & same \\
$\F:\Sw\to\Sw$ continuous & Theorem~\ref{prop:F-S} & same \\
$\ip{\F T}{\varphi}=\ip{T}{\F\varphi}$ & \eqref{eq:dualdef} & Proposition~\ref{prop:dual-consistency} \\
Plancherel $\|f\|_{2}^{2}=(2\pi)^{-1}\|\fh\|_{2}^{2}$ & Theorem~\ref{thm:plancherel} & same \\
\bottomrule
\end{tabular}
\end{center}

The second column records where a statement is first \emph{computed}, the third where it is
\emph{derived from the definitions}. The gap between them is the content of
Chapters~\ref{ch:limits}--\ref{ch:tespered}: the early chapters compute with rules that are only
justified later, and the later ones supply the justification without changing any of the numbers.
